\documentclass[11pt]{article}

\usepackage[utf8]{inputenc}
\usepackage[T1]{fontenc}
\usepackage[margin=0.8in]{geometry}
\usepackage{amsmath}
\usepackage{booktabs}
\usepackage{enumitem}
\usepackage{array}
\usepackage{hyperref}

\setlist[itemize]{leftmargin=1.3em, noitemsep, topsep=3pt}
\setlist[enumerate]{leftmargin=1.5em, noitemsep, topsep=3pt}

\title{\textbf{System Design Revision Notes}}
\author{}
\date{}

\begin{document}

\maketitle

\section{Distributed System}

\subsection*{Definition}

A distributed system is a collection of independent components located on different machines that share messages with each other in order to achieve common goals.

\section{Performance}

\subsection*{Definition}

Performance describes how a system behaves under a specific workload. It is not one number. It is measured using several metrics, such as latency, throughput, error rate, resource utilization, queue length, uptime, and recovery time.

\[
\boxed{\text{Performance} = \text{System behavior under a specific workload}}
\]

\subsection*{Example}

Suppose a payment system receives 1,000 requests per second. Under this load, it completes 950 requests per second, average latency is 300 ms, error rate is 2\%, and CPU utilization is 75\%. Together, these values describe the system's performance.

\subsection*{Important Formulas}

\[
\text{Error Rate}
=
\frac{\text{Failed Requests}}{\text{Total Requests}}
\times 100
\]

\[
\text{Uptime}
=
\frac{\text{Total Time} - \text{Downtime}}{\text{Total Time}}
\times 100
\]

\[
\text{MTBF}
=
\frac{\text{Total Operating Time}}{\text{Number of Failures}}
\]

\[
\text{MTTR}
=
\frac{\text{Total Recovery Time}}{\text{Number of Failures}}
\]

\subsection*{Solved Problem}

\textbf{Problem:} A system processes 3,600 requests in 12 hours. During this period, 72 requests fail, the system experiences 3 outages, and each outage lasts 15 minutes.

\textbf{Find:} Throughput, error rate, MTBF, MTTR, and uptime.

\textbf{Solution:}

\[
\text{Throughput}
=
\frac{3600}{12}
=
300 \text{ requests/hour}
\]

\[
\text{Error Rate}
=
\frac{72}{3600}
\times 100
=
2\%
\]

\[
\text{MTBF}
=
\frac{12}{3}
=
4 \text{ hours}
\]

\[
\text{MTTR}
=
15 \text{ minutes}
\]

Total downtime is:

\[
3 \times 15 = 45 \text{ minutes}
\]

Total time is:

\[
12 \times 60 = 720 \text{ minutes}
\]

\[
\text{Uptime}
=
\frac{720 - 45}{720}
\times 100
=
93.75\%
\]

\textbf{Final Answer:} Throughput is 300 requests/hour, error rate is 2\%, MTBF is 4 hours, MTTR is 15 minutes, and uptime is 93.75\%.

\section{Scalability}

\subsection*{Definition}

Scalability is the ability of a system to handle increasing demand by adding resources while maintaining acceptable performance.

\[
\boxed{\text{Scalability} = \text{Ability to handle growth by adding resources}}
\]

\subsection*{Types of Scaling}

\textbf{Vertical scaling} means making one machine more powerful, such as adding CPU, memory, or faster storage.

\textbf{Horizontal scaling} means adding more machines or server instances.

\subsection*{Important Concepts and Formulas}

\textbf{Capacity} is the maximum workload a system can sustainably handle while meeting latency and error-rate requirements.

\[
\text{Resource Multiplier}
=
\frac{\text{New Resources}}{\text{Original Resources}}
\]

\[
\text{Capacity Multiplier}
=
\frac{\text{New Capacity}}{\text{Original Capacity}}
\]

\[
\text{Capacity Increase}
=
\frac{\text{New Capacity} - \text{Original Capacity}}
{\text{Original Capacity}}
\times 100
\]

\[
\text{Scaling Efficiency}
=
\frac{\text{Capacity Multiplier}}
{\text{Resource Multiplier}}
\times 100
\]

\textbf{Elasticity} measures how quickly a system adds or removes resources when demand changes.

\[
\boxed{\text{Elasticity} = \text{Speed of adapting resources to changing demand}}
\]

\textbf{Cost efficiency} measures how much useful capacity is obtained for the money spent.

\[
\text{Capacity per Dollar}
=
\frac{\text{System Capacity}}{\text{Hourly Cost}}
\]

\subsection*{Solved Problem}

\textbf{Problem:} A system uses 4 servers and supports 1,200 requests per second. The number of servers is increased to 8, and the new system supports 2,040 requests per second.

\textbf{Find:} Resource multiplier, capacity multiplier, capacity increase, and scaling efficiency.

\textbf{Solution:}

\[
\text{Resource Multiplier}
=
\frac{8}{4}
=
2
\]

\[
\text{Capacity Multiplier}
=
\frac{2040}{1200}
=
1.7
\]

\[
\text{Capacity Increase}
=
\frac{2040 - 1200}{1200}
\times 100
=
70\%
\]

\[
\text{Scaling Efficiency}
=
\frac{1.7}{2}
\times 100
=
85\%
\]

\textbf{Final Answer:} Resources increased by $2\times$, capacity increased by $1.7\times$, capacity increased by 70\%, and scaling efficiency is 85\%.

\textbf{Interpretation:} The system scales horizontally, but it does not scale perfectly because doubling the servers does not double the capacity.

\subsection*{Elasticity Example}

A streaming system has 5 servers, and each server supports 200 streams.

\[
\text{Original Capacity}
=
5 \times 200
=
1000 \text{ streams}
\]

Demand increases to 1,600 streams.

\[
\text{Required Servers}
=
\frac{1600}{200}
=
8
\]

\[
\text{Additional Servers Required}
=
8 - 5
=
3
\]

If autoscaling takes 5 minutes, then the scale-up response time is 5 minutes. Before scaling finishes, the capacity shortfall is:

\[
1600 - 1000 = 600 \text{ streams}
\]

\section{Latency}

\subsection*{Definition}

Latency is the total time between sending a request and receiving its response.

\[
\boxed{\text{Latency} = \text{Time taken by one request}}
\]

\[
\text{Latency}
=
\text{Response Time}
-
\text{Request Start Time}
\]

For a payment request, latency starts when the user clicks the Pay button and ends when the user receives confirmation.

\subsection*{Components of Latency}

\[
\begin{aligned}
\text{Total Latency}
={}&
\text{Network Delay}
+
\text{Queueing Delay} \\
&+
\text{Processing Delay}
+
\text{Database Delay} \\
&+
\text{External Service Delay}
\end{aligned}
\]

\begin{itemize}
    \item \textbf{Network delay:} Time required for data to travel.
    \item \textbf{Queueing delay:} Time spent waiting before processing.
    \item \textbf{Processing delay:} Time required by the server.
    \item \textbf{Database delay:} Time required to read or write data.
    \item \textbf{External service delay:} Time spent waiting for another system.
\end{itemize}

\subsection*{Queueing Example}

Suppose a server processes 100 requests per second, and 700 requests arrive simultaneously. The requests are processed in groups of 100.

\begin{center}
\begin{tabular}{ccc}
\toprule
\textbf{Group} & \textbf{Requests} & \textbf{Approximate Latency} \\
\midrule
1 & 1--100 & 1 second \\
2 & 101--200 & 2 seconds \\
3 & 201--300 & 3 seconds \\
4 & 301--400 & 4 seconds \\
5 & 401--500 & 5 seconds \\
6 & 501--600 & 6 seconds \\
7 & 601--700 & 7 seconds \\
\bottomrule
\end{tabular}
\end{center}

The 700th request has an approximate latency of 7 seconds. Its processing time may still be one second, but it experiences additional queueing delay.

\subsection*{Timeout Example}

Suppose capacity is 100 requests per second, 700 requests arrive, and the timeout is 5 seconds.

Requests completed before timeout:

\[
100 \times 5 = 500
\]

Timed-out requests:

\[
700 - 500 = 200
\]

Timeout error rate:

\[
\frac{200}{700}
\times 100
=
28.57\%
\]

\textbf{Final Answer:} 500 requests complete, 200 time out, maximum successful latency is 5 seconds, and timeout error rate is 28.57\%.

\subsection*{Percentile Latency: P50, P95, and P99}

\subsection*{Definition}

Percentile latency shows latency distribution across many requests.

\begin{itemize}
    \item \textbf{P50 latency:} Median latency; 50\% of requests finish within this time.
    \item \textbf{P95 latency:} 95\% of requests finish within this time; the slowest 5\% take longer.
    \item \textbf{P99 latency:} 99\% of requests finish within this time; the slowest 1\% take longer.
\end{itemize}

\[
\boxed{
\text{P50 shows typical latency; P95 and P99 expose slow-user experience.}
}
\]

\subsection*{Example}

Suppose 10 payment requests have these sorted latencies:

\[
100,\ 120,\ 150,\ 180,\ 200,\ 250,\ 300,\ 400,\ 800,\ 1200 \text{ ms}
\]

Using the nearest-rank method:

\[
\text{P50 position} = 0.50 \times 10 = 5
\]

\[
\text{P50 latency} = 200 \text{ ms}
\]

\[
\text{P95 position} = 0.95 \times 10 = 9.5 \approx 10
\]

\[
\text{P95 latency} = 1200 \text{ ms}
\]

\subsection*{Interpretation}

Most requests are fast, but a few are very slow. Average latency may hide this problem, while P95 and P99 reveal it.

\[
\boxed{
\text{Average latency hides outliers; percentile latency exposes them.}
}

\section{Throughput}

\subsection*{Definition}

Throughput is the amount of work completed by a system during a given period.

\[
\boxed{\text{Throughput} = \frac{\text{Completed Work}}{\text{Time}}}
\]

\[
\text{Throughput}
=
\frac{\text{Completed Requests}}{\text{Time}}
\]

Common units include requests per second, transactions per second, and messages per second.

\subsection*{Important Distinctions}

\[
\text{Incoming Rate}
=
\frac{\text{Requests Received}}{\text{Time}}
\]

\[
\text{Completed Throughput}
=
\frac{\text{Completed Requests}}{\text{Time}}
\]

\[
\text{Goodput}
=
\frac{\text{Successful Requests}}{\text{Time}}
\]

\textbf{Incoming request rate} counts requests arriving at the system. \textbf{Completed throughput} counts requests that finish processing. \textbf{Goodput} counts only successful and useful requests.

\subsection*{Example}

During one second, 100 requests arrive, 95 succeed, 3 finish with errors, and 2 remain queued.

\[
\text{Incoming Rate}
=
100 \text{ requests/second}
\]

\[
\text{Completed Throughput}
=
95 + 3
=
98 \text{ requests/second}
\]

\[
\text{Goodput}
=
95 \text{ requests/second}
\]

The queued requests are not counted because they have not finished.

\subsection*{Solved Problem}

A system receives 700 requests and processes 500 requests within 5 seconds. Of the 500 completed requests, 480 succeed and 20 return errors. Another 200 requests time out.

Completed throughput:

\[
\frac{500}{5}
=
100 \text{ requests/second}
\]

Successful throughput:

\[
\frac{480}{5}
=
96 \text{ requests/second}
\]

Overall unsuccessful requests:

\[
20 + 200 = 220
\]

Overall unsuccessful rate:

\[
\frac{220}{700}
\times 100
=
31.43\%
\]

\textbf{Final Answer:} Completed throughput is 100 requests/second, successful throughput is 96 requests/second, and the overall unsuccessful rate is 31.43\%.

\section{Consistency}

\subsection*{Definition}

\begin{itemize}
    \item Every read returns the latest successful write.
    \item All replicas provide the same view of the data.
    \item Users do not observe stale data.
\end{itemize}

\[
\boxed{\text{Consistency} = \text{Latest data is always returned}}
\]

\subsection*{Example}

\begin{itemize}
    \item A user transfers \$500 from their bank account.
    \item The balance changes from \$2,000 to \$1,500.
    \item If every ATM, mobile app, and bank server shows \$1,500, the system is consistent.
    \item If one ATM still shows \$2,000, the system is temporarily inconsistent.
\end{itemize}

\section{Availability}

\subsection*{Definition}

\begin{itemize}
    \item Every request receives a response.
    \item The response may contain stale data.
    \item Availability focuses on responding, not necessarily with the latest data.
\end{itemize}

\[
\boxed{\text{Availability} = \text{Every request gets a response}}
\]

\subsection*{Example}

\begin{itemize}
    \item East Coast and West Coast Instagram data centers lose communication.
    \item Both continue serving their local users.
    \item Users receive responses even though the data may differ.
\end{itemize}

\section{Partition Tolerance}

\subsection*{Definition}

\begin{itemize}
    \item The system continues operating despite a network partition.
    \item A network partition means two servers or data centers cannot communicate.
    \item It is unrelated to Kafka partitions, database partitioning, or sharding.
\end{itemize}

\[
\boxed{\text{Partition Tolerance} = \text{System survives network failures}}
\]

\subsection*{Example}

\begin{itemize}
    \item East Coast and West Coast data centers become disconnected.
    \item Both continue serving local users.
    \item Updates cannot be replicated until communication is restored.
\end{itemize}

\section{CAP Theorem}

\subsection*{Definition}

During a network partition, a distributed system can guarantee only one of:

\[
\boxed{\text{Consistency + Partition Tolerance (CP)}}
\]

or

\[
\boxed{\text{Availability + Partition Tolerance (AP)}}
\]

\[
\boxed{\text{CAP: During a network partition, choose CP or AP}}
\]

\subsection*{CP System}

\begin{itemize}
    \item Prioritizes consistency.
    \item May reject or delay requests.
    \item Never returns stale data.
\end{itemize}

\subsection*{AP System}

\begin{itemize}
    \item Prioritizes availability.
    \item Always responds to requests.
    \item May temporarily return stale data.
\end{itemize}

\subsection*{Instagram Example}

\begin{itemize}
    \item A user deletes a post on the West Coast.
    \item Due to a network partition, the East Coast has not received the update.
    \item \textbf{AP:} East Coast still shows the old post.
    \item \textbf{CP:} East Coast blocks or delays the request until it confirms the latest data.
\end{itemize}

\section{Weak Consistency}

\subsection*{Definition}

Weak consistency does not guarantee that a read immediately following a write will return the most recently updated value. Different users or replicas may temporarily observe different versions of the same data.

\noindent
\fbox{\parbox{0.95\linewidth}{
\textbf{Quick Revision:} Weak consistency allows stale reads and provides no specific guarantee about when the latest value will become visible.
}}

\begin{itemize}
    \item A write may not be immediately visible to all users.
    \item Subsequent reads may return old or stale data.
    \item Different replicas may temporarily contain different values.
    \item Weak consistency alone does not guarantee when replicas will converge.
\end{itemize}

\subsection*{Real-World Example}

Suppose an e-commerce platform changes a product price from \$100 to \$120. Immediately after the update, one user may see \$120 while another user connected to a different replica may still see \$100.

\section{Eventual Consistency}

\subsection*{Definition}

Eventual consistency is a form of weak consistency in which replicas may temporarily contain different values, but they will eventually converge to the same value if no new writes occur.

\noindent
\fbox{\parbox{0.95\linewidth}{
\textbf{Quick Revision:} Eventual consistency allows temporary stale reads but guarantees that all replicas will eventually converge.
}}

\begin{itemize}
    \item Updates are commonly propagated asynchronously.
    \item Reads may temporarily return stale data.
    \item Replicas eventually become consistent if no new updates occur.
    \item It provides higher availability and lower write latency than strong consistency in many distributed systems.
\end{itemize}

\subsection*{Real-World Example}

A user updates their profile picture on a social-media platform. Some users may temporarily see the old picture because they are reading from replicas that have not yet received the update. After replication completes, all users see the new picture.

\subsection*{Important Distinction}

\begin{itemize}
    \item \textbf{Weak consistency} does not necessarily guarantee convergence.
    \item \textbf{Eventual consistency} guarantees convergence when no additional writes occur.
\end{itemize}

\section{Strong Consistency}

\subsection*{Definition}

Strong consistency guarantees that once a write operation successfully completes, every subsequent read returns the latest completed value.

\noindent
\fbox{\parbox{0.95\linewidth}{
\textbf{Quick Revision:} Strong consistency ensures that every read after a completed write observes the latest value.
}}

\begin{itemize}
    \item Stale reads are not allowed after a write has completed.
    \item The system behaves as if there is a single up-to-date copy of the data.
    \item Writes may require coordination between a leader or a quorum of replicas.
    \item Strong consistency may increase latency and reduce availability during network failures.
\end{itemize}

\subsection*{Real-World Example}

A bank account has a balance of \$1,000. After a successful \$200 withdrawal, every subsequent balance check must return \$800. Returning the previous balance of \$1,000 would violate strong consistency.

\subsection*{Important Distinction}

Strong consistency does not necessarily mean that every physical replica is updated at exactly the same instant. It means that the system prevents clients from observing stale data after the write has been acknowledged.

% Replacement/additional rows for the Final Comparison table

\section{Availability}

\subsection*{Definition}
Availability is the percentage of time a system remains operational and accessible during a specified period.

\begin{center}
\fbox{\parbox{0.92\linewidth}{\centering
Availability measures how often a system is usable.
}}
\end{center}

\[
\text{Availability}
=
\frac{\text{Uptime}}
{\text{Uptime}+\text{Downtime}}
\times 100
\]

\begin{itemize}
    \item Uptime is the duration for which the system operates correctly.
    \item Downtime is the duration for which the system is unavailable.
    \item Replication, redundancy, and failover are commonly used to improve availability.
\end{itemize}

\subsection*{Real-World Example}
An e-commerce system uses multiple application servers and a backup database so that users can continue placing orders even when one component fails.


\section{Availability Nines}

\subsection*{Definition}
Availability is commonly expressed using ``nines.'' The number of consecutive \(9\)s represents the expected availability level of the system.

\begin{center}
\fbox{\parbox{0.92\linewidth}{\centering
More nines mean higher availability and less permitted downtime.
}}
\end{center}

\begin{itemize}
    \item \(99.9\%\) availability is called three nines.
    \item \(99.99\%\) availability is called four nines.
    \item \(99.999\%\) availability is called five nines.
    \item Achieving additional nines usually requires more redundancy, complexity, and cost.
\end{itemize}

\[
\text{Allowed Downtime}
=
\text{Total Time}
\times
(1-\text{Availability})
\]

\subsection*{Approximate Annual Downtime}

\begin{center}
\begin{tabular}{|c|c|c|}
\hline
\textbf{Availability} & \textbf{Name} & \textbf{Annual Downtime} \\
\hline
\(99\%\) & Two nines & Approximately \(3.65\) days \\
\hline
\(99.9\%\) & Three nines & Approximately \(8.76\) hours \\
\hline
\(99.99\%\) & Four nines & Approximately \(52.56\) minutes \\
\hline
\(99.999\%\) & Five nines & Approximately \(5.26\) minutes \\
\hline
\end{tabular}
\end{center}

\subsection*{Solved Problem}

\textbf{Problem:} Calculate the maximum annual downtime for a system with \(99.99\%\) availability.

\textbf{Given:}
\begin{itemize}
    \item Availability \(= 99.99\% = 0.9999\)
    \item Total minutes in one year \(= 365 \times 24 \times 60 = 525{,}600\)
\end{itemize}

\textbf{Find:} Maximum annual downtime.

\textbf{Solution:}
\[
\text{Downtime}
=
525{,}600 \times (1-0.9999)
\]

\[
\text{Downtime}
=
525{,}600 \times 0.0001
=
52.56 \text{ minutes}
\]

\textbf{Final Answer:}
\[
\boxed{52.56 \text{ minutes per year}}
\]


\section{Replication}

\subsection*{Definition}
Replication is the process of maintaining copies of the same data on multiple servers to improve availability and fault tolerance.

\begin{center}
\fbox{\parbox{0.92\linewidth}{\centering
Replication keeps multiple copies of data so that another server can continue serving requests after a failure.
}}
\end{center}

\begin{itemize}
    \item Replicas reduce the risk of a single database server becoming a single point of failure.
    \item Replication may be synchronous or asynchronous.
    \item Replication alone does not guarantee automatic recovery; a failover mechanism may also be required.
\end{itemize}

\subsection*{Real-World Example}
An online banking system stores account data on a primary database and one or more replica databases. If the primary database fails, a replica can be promoted to replace it.


\section{Primary-Primary Replication}

\subsection*{Definition}
Primary-primary replication is a replication pattern in which multiple database servers accept both read and write operations and replicate changes to each other.

\begin{center}
\fbox{\parbox{0.92\linewidth}{\centering
Primary-primary replication improves availability and write capacity but introduces conflict-resolution complexity.
}}
\end{center}

\begin{itemize}
    \item All primary nodes can serve read and write requests.
    \item The system can continue accepting writes when one primary node fails.
    \item Write conflicts may occur when different nodes update the same data simultaneously.
    \item Conflict detection and resolution are required.
\end{itemize}

\subsection*{Advantages}
\begin{itemize}
    \item High availability for both reads and writes.
    \item Better utilization of all database nodes.
    \item Can support users from multiple geographic regions.
\end{itemize}

\subsection*{Disadvantages}
\begin{itemize}
    \item Complex conflict resolution.
    \item Greater risk of inconsistent data.
    \item Replication and operational management are more difficult.
\end{itemize}

\subsection*{Real-World Example}
A globally distributed application allows users in the United States and Europe to write to nearby database nodes. Both nodes replicate updates to each other, and the system must resolve conflicting updates to the same record.


\section{Primary-Replica Replication}

\subsection*{Definition}
Primary-replica replication is a pattern in which one primary server accepts write operations, while replica servers copy its data and typically serve read operations.

\begin{center}
\fbox{\parbox{0.92\linewidth}{\centering
Primary-replica replication simplifies writes by allowing only one primary writer.
}}
\end{center}

\begin{itemize}
    \item The primary server handles writes and may also handle reads.
    \item Replica servers usually handle read requests.
    \item Replicas continuously copy changes from the primary.
    \item If the primary fails, a replica may be promoted to become the new primary.
\end{itemize}

\subsection*{Advantages}
\begin{itemize}
    \item Simpler consistency management because only one server accepts writes.
    \item Read traffic can be distributed across multiple replicas.
    \item Easier conflict handling than primary-primary replication.
\end{itemize}

\subsection*{Disadvantages}
\begin{itemize}
    \item The primary may become a write bottleneck.
    \item Replicas may return stale data because of replication delay.
    \item Failover may cause temporary downtime.
\end{itemize}

\subsection*{Real-World Example}
An e-commerce database sends all product updates to the primary server, while read replicas serve product-search and product-detail requests.


\section{Failover}

\subsection*{Definition}
Failover is the process of transferring traffic or responsibilities from a failed component to a healthy replacement component.

\begin{center}
\fbox{\parbox{0.92\linewidth}{\centering
Failover keeps a system available by replacing a failed component with a healthy one.
}}
\end{center}

\begin{itemize}
    \item Failover may be automatic or manual.
    \item Health checks or heartbeat messages are used to detect failures.
    \item The replacement component must have the required data and configuration.
    \item Failover time directly affects system downtime.
\end{itemize}

\subsection*{Real-World Example}
If the primary database stops sending heartbeat signals, the system promotes a replica database and redirects application traffic to it.


\section{Active-Active Failover}

\subsection*{Definition}
Active-active is a failover pattern in which multiple components actively serve traffic at the same time.

\begin{center}
\fbox{\parbox{0.92\linewidth}{\centering
In active-active failover, all healthy components serve traffic and share the workload.
}}
\end{center}

\begin{itemize}
    \item Traffic is distributed across multiple active components.
    \item If one component fails, the remaining components absorb its traffic.
    \item A load balancer commonly routes requests to healthy components.
    \item Stateless services are well suited for this pattern.
\end{itemize}

\subsection*{Advantages}
\begin{itemize}
    \item High resource utilization.
    \item Minimal interruption during failure.
    \item Supports horizontal scaling.
\end{itemize}

\subsection*{Disadvantages}
\begin{itemize}
    \item Data synchronization can be complex for stateful components.
    \item Remaining nodes must have enough capacity to handle additional traffic.
    \item Split-brain and conflicting writes may occur in distributed databases.
\end{itemize}

\subsection*{Real-World Example}
Several stateless web servers run behind a load balancer. All servers process requests, and if one server fails, the load balancer sends traffic to the remaining servers.


\section{Active-Passive Failover}

\subsection*{Definition}
Active-passive is a failover pattern in which one component actively serves traffic while another component remains on standby.

\begin{center}
\fbox{\parbox{0.92\linewidth}{\centering
In active-passive failover, the standby component takes over only after the active component fails.
}}
\end{center}

\begin{itemize}
    \item The active component handles normal traffic.
    \item The passive component remains ready to take over.
    \item Heartbeats or health checks are used to detect active-component failure.
    \item During failover, the passive component may assume the active component's network address or role.
\end{itemize}

\subsection*{Advantages}
\begin{itemize}
    \item Simpler consistency management.
    \item Lower risk of simultaneous conflicting writes.
    \item Easier to operate than active-active for stateful databases.
\end{itemize}

\subsection*{Disadvantages}
\begin{itemize}
    \item The standby component is underutilized during normal operation.
    \item Failover may introduce a short interruption.
    \item The standby component may contain slightly stale data.
\end{itemize}

\subsection*{Real-World Example}
A primary database processes all reads and writes, while a standby database continuously receives replicated updates. If the primary fails, the standby is promoted and begins serving application traffic.


\section{Sequential Availability}

\subsection*{Definition}
Components are arranged sequentially when every component must be available for the complete system operation to succeed.

\begin{center}
\fbox{\parbox{0.92\linewidth}{\centering
In a sequential system, failure of any required component can make the entire system unavailable.
}}
\end{center}

For independent sequential components:

\[
A_{\text{system}}
=
A_1 \times A_2 \times \cdots \times A_n
\]

\begin{itemize}
    \item Every component is part of the critical execution path.
    \item Overall availability is lower than the availability of each individual component.
    \item Adding more mandatory dependencies decreases total system availability.
\end{itemize}

\subsection*{Real-World Example}
An e-commerce checkout request sequentially depends on the order service, payment service, inventory service, and database. If any required component is unavailable, checkout fails.

\subsection*{Solved Problem}

\textbf{Problem:} Calculate the availability of two sequential components.

\textbf{Given:}
\begin{itemize}
    \item Component \(A\) availability \(= 99.99\% = 0.9999\)
    \item Component \(B\) availability \(= 99.99\% = 0.9999\)
\end{itemize}

\textbf{Find:} Total system availability.

\textbf{Solution:}
\[
A_{\text{system}}
=
0.9999 \times 0.9999
=
0.99980001
\]

\[
A_{\text{system}}
=
99.980001\%
\]

\textbf{Final Answer:}
\[
\boxed{A_{\text{system}} \approx 99.98\%}
\]


\section{Parallel Availability}

\subsection*{Definition}
Components are arranged in parallel when the system can continue operating as long as at least one redundant component remains available.

\begin{center}
\fbox{\parbox{0.92\linewidth}{\centering
Parallel redundancy increases availability because the system fails only when all redundant components fail.
}}
\end{center}

For two independent parallel components:

\[
A_{\text{system}}
=
1-(1-A_1)(1-A_2)
\]

For \(n\) identical parallel components:

\[
A_{\text{system}}
=
1-(1-A)^n
\]

\begin{itemize}
    \item Parallel components provide alternative execution paths.
    \item Failure of one component does not necessarily stop the system.
    \item Load balancing or failover logic is required to use healthy components.
\end{itemize}

\subsection*{Real-World Example}
Two application servers operate behind a load balancer. If one server fails, the other server continues processing requests.

\subsection*{Solved Problem}

\textbf{Problem:} Calculate the availability of two parallel servers.

\textbf{Given:}
\begin{itemize}
    \item Server \(A\) availability \(= 99.99\% = 0.9999\)
    \item Server \(B\) availability \(= 99.99\% = 0.9999\)
\end{itemize}

\textbf{Find:} Total system availability.

\textbf{Solution:}
\[
A_{\text{system}}
=
1-(1-0.9999)(1-0.9999)
\]

\[
A_{\text{system}}
=
1-(0.0001)(0.0001)
\]

\[
A_{\text{system}}
=
1-0.00000001
=
0.99999999
\]

\[
A_{\text{system}}
=
99.999999\%
\]

\textbf{Final Answer:}
\[
\boxed{A_{\text{system}} = 99.999999\%}
\]


\section{Choosing an Availability Pattern}

\subsection*{Definition}
An availability pattern should be selected according to the component's state, consistency requirements, traffic load, recovery-time requirement, and infrastructure cost.

\begin{center}
\fbox{\parbox{0.92\linewidth}{\centering
Use active-active for easily distributed workloads and active-passive when simpler state management is more important.
}}
\end{center}

\begin{itemize}
    \item Stateless application servers commonly use active-active deployment.
    \item Stateful databases commonly use active-passive or primary-replica deployment.
    \item Primary-primary replication is useful when writes must remain available in multiple locations.
    \item Primary-replica replication is preferred when simpler consistency and read scaling are required.
\end{itemize}

\subsection*{Real-World Example}
An e-commerce system may use active-active web servers behind a load balancer and an active-passive primary-replica database. The web servers provide scaling and rapid traffic redistribution, while the database avoids simultaneous write conflicts.

\section{Background Tasks}

\subsection*{Definition}

A background task is work that runs outside the main user-facing request-response flow. It is used when work should continue without blocking the user interface or making the user wait for a long operation to finish.

\[
\boxed{\text{A background task runs separately from the main user request.}}
\]

\begin{itemize}
    \item It usually runs in a separate worker process or service.
    \item It is useful for slow, delayed, or non-critical work.
    \item It improves perceived latency for the user.
    \item Common examples include sending emails, processing files, generating reports, and running data pipelines.
    \item Background tasks can be event-driven, scheduled, or polling-based.
\end{itemize}

\subsection*{Real-World System-Design Example}

In a meeting-summary system, the user uploads a transcript through the main application. Instead of generating the full summary during the upload request, the system stores the transcript and lets background workers process it later.


\section{Event-Driven Background Tasks}

\subsection*{Definition}

An event-driven background task is a background task that starts when a specific event occurs. The event may be a user action, database change, message arrival, file upload, or API call.

\[
\boxed{\text{Event-driven tasks run because an event occurred.}}
\]

\begin{itemize}
    \item A producer creates an event, such as \texttt{OrderPlaced} or \texttt{FileUploaded}.
    \item The event is published to a queue, message broker, or pub-sub system.
    \item A background worker consumes the event and processes it.
    \item Event-driven tasks are useful when the system must react quickly to new work.
    \item They usually require retries, dead-letter queues, and idempotent processing.
    \item They are more responsive than scheduled or polling-based tasks.
\end{itemize}

\subsection*{Real-World System-Design Example}

In a food-delivery application, when a user places an order, the system publishes an \texttt{OrderPlaced} event. Background workers then notify the restaurant, process the payment, and start driver assignment.

\subsection*{Important Distinction}

An event-driven task is different from a polling-based task because the worker does not repeatedly check whether work exists. Instead, work is pushed to the worker through an event, queue, or message broker.


\section{Scheduled Background Tasks}

\subsection*{Definition}

A scheduled background task is a background task that runs at a predefined time or fixed interval. It is triggered by a clock, cron job, or scheduling service.

\[
\boxed{\text{Scheduled tasks run because the configured time has arrived.}}
\]

\begin{itemize}
    \item The task runs at a fixed time, such as every day at 2:00 AM.
    \item It may also run at a fixed interval, such as every hour or every 10 minutes.
    \item It runs whether or not new work has arrived.
    \item It is useful for predictable maintenance work.
    \item Common examples include backups, cleanup jobs, billing reports, and daily summaries.
    \item It is simple and predictable, but may waste resources if there is no useful work.
\end{itemize}

\subsection*{Real-World System-Design Example}

In a food-delivery platform, a scheduled job runs every night at 2:00 AM to calculate restaurant earnings and generate daily settlement reports.

\subsection*{Important Distinction}

A scheduled task is time-driven. The main reason it runs is not because new work appeared, but because the configured time or interval was reached.


\section{Polling-Based Background Tasks}

\subsection*{Definition}

A polling-based background task is a background task that repeatedly checks a database, queue, API, or system condition to see whether there is work to do.

\[
\boxed{\text{Polling-based tasks repeatedly check for work and act only when work exists.}}
\]

\begin{itemize}
    \item The worker runs continuously in the background.
    \item It checks for eligible work at regular intervals.
    \item If work exists, it processes the work.
    \item If no work exists, it sleeps and checks again later.
    \item It is useful when the system being monitored cannot directly publish events.
    \item Short polling intervals reduce delay but increase resource usage.
    \item Long polling intervals reduce resource usage but increase detection delay.
\end{itemize}

\subsection*{Real-World System-Design Example}

In a meeting-summary system, a database worker checks every few seconds for meetings that have uploaded transcripts but missing summary artifacts. If it finds such meetings, it creates workflow tasks for the next processing stage.

\subsection*{Important Distinction}

Polling-based tasks are not purely scheduled tasks. A scheduled task runs because the time has arrived. A polling task wakes up repeatedly to check whether a condition has become true, and only performs meaningful work if the condition exists.


\section{Comparison of Background Task Trigger Types}

\subsection*{Definition}

Background tasks can be classified by what triggers them: an event, a schedule, or repeated polling. The trigger model determines the task's latency, complexity, resource usage, and reliability requirements.

\[
\boxed{\text{The trigger model decides when a background task starts working.}}
\]

\begin{itemize}
    \item Event-driven tasks are best when fast reaction is required.
    \item Scheduled tasks are best when work must happen at predictable times.
    \item Polling-based tasks are best when the system must detect state changes without direct event notifications.
    \item In interviews, explain both the trigger mechanism and the trade-offs.
\end{itemize}

\subsection*{Real-World System-Design Example}

In a food-delivery system, order placement can use event-driven workers, nightly settlement reports can use scheduled jobs, and stale-delivery detection can use polling workers.

\section{Domain Name System (DNS)}

\subsection*{Definition}
The Domain Name System (DNS) is a distributed lookup system that translates human-readable domain names, such as \texttt{google.com}, into IP addresses that computers use to connect to servers.

\[
\boxed{\text{DNS does not fetch the web page; it only finds the IP address for a domain name.}}
\]

\begin{itemize}
    \item A browser uses DNS when it needs to convert a domain name into an IP address.
    \item DNS works like a phone book for the internet.
    \item The browser first checks local cache before asking external DNS systems.
    \item After receiving the IP address, the browser connects directly to the web server.
\end{itemize}

\subsection*{Real-World System-Design Example}
When a user types \texttt{www.google.com} in the browser:
\begin{itemize}
    \item The browser checks whether it already knows the IP address.
    \item If not, it asks the DNS resolver.
    \item The DNS system returns the IP address for \texttt{www.google.com}.
    \item The browser then uses that IP address to connect to Google's web server.
\end{itemize}

\subsection*{Important Distinction}
\begin{itemize}
    \item DNS resolves a domain name to an IP address.
    \item The web server returns the actual web page content.
    \item DNS is part of connection setup, not page rendering.
\end{itemize}


\section{DNS Lookup Flow}

\subsection*{Definition}
DNS lookup flow is the step-by-step process used to find the IP address of a domain name when it is not already available in cache.

\[
\boxed{\text{DNS lookup follows a hierarchy: browser cache, resolver, root server, TLD server, and authoritative server.}}
\]

\begin{itemize}
    \item The browser first checks its own DNS cache.
    \item If the answer is not found, the request goes to the operating system or router.
    \item The request is usually forwarded to a recursive DNS resolver.
    \item If the resolver does not know the answer, it queries DNS servers in a hierarchy.
    \item The final IP address comes from the authoritative DNS server.
\end{itemize}

\subsection*{Real-World System-Design Example}
For \texttt{www.example.com}:
\begin{itemize}
    \item The browser checks whether \texttt{www.example.com} is already cached.
    \item If not cached, the request goes to the DNS resolver.
    \item The resolver may ask a root server where to find \texttt{.com}.
    \item The \texttt{.com} TLD server points to the authoritative server for \texttt{example.com}.
    \item The authoritative server returns the IP address.
    \item The browser connects to that IP address.
\end{itemize}

\subsection*{Important Distinction}
\begin{itemize}
    \item A root server does not usually return the final IP address.
    \item A TLD server usually points to the domain's authoritative DNS server.
    \item The authoritative DNS server provides the final DNS answer.
\end{itemize}


\section{Devices Involved in DNS Resolution}

\subsection*{Definition}
DNS resolution involves multiple devices and servers that cooperate to convert a domain name into an IP address.

\[
\boxed{\text{DNS resolution starts from the user's device and may pass through routers, resolvers, and DNS hierarchy servers.}}
\]

\begin{itemize}
    \item The user's computer or phone starts the request.
    \item The browser checks its DNS cache.
    \item The operating system may also check its local DNS cache.
    \item The home or office router may forward the DNS request.
    \item A recursive DNS resolver performs the lookup on behalf of the user.
    \item Root, TLD, and authoritative DNS servers may be contacted during the lookup.
\end{itemize}

\subsection*{Real-World System-Design Example}
A laptop connected to home Wi-Fi opens \texttt{amazon.com}:
\begin{itemize}
    \item The laptop browser checks its cache.
    \item The laptop sends the DNS request through the Wi-Fi router.
    \item The router forwards the request to a DNS resolver.
    \item The resolver finds the correct IP address.
    \item The IP address is returned to the laptop.
    \item The laptop browser connects to Amazon's server using that IP address.
\end{itemize}

\subsection*{Important Distinction}
\begin{itemize}
    \item The router usually forwards DNS traffic; it is not the full DNS system.
    \item The recursive resolver performs the lookup work for the user.
    \item The final website connection happens after DNS resolution is complete.
\end{itemize}

\section{Content Delivery Network (CDN)}

\subsection*{Definition}
A Content Delivery Network (CDN) is a distributed network of edge servers that stores and serves content closer to users instead of sending every request directly to the origin server.

\[
\boxed{\text{A CDN reduces latency by serving cached content from nearby edge servers.}}
\]

\begin{itemize}
    \item The origin server is the main backend or storage system where the original content lives.
    \item CDN edge servers are placed closer to users in different geographic locations.
    \item If requested content is available at the edge server, the CDN returns it directly.
    \item If the content is not available, the CDN fetches it from the origin server.
    \item CDNs reduce latency, improve availability, and reduce load on the origin server.
\end{itemize}

\subsection*{Example}
When a user in New York watches a Netflix video, the video is usually served from a nearby CDN edge server instead of Netflix's main origin servers. This reduces buffering and improves playback speed.

\subsection*{Basic Request Flow}
\begin{itemize}
    \item User requests content, such as an image, video, or web page.
    \item DNS resolves the request to a nearby CDN edge server.
    \item CDN checks whether the content exists in its cache.
    \item If cache hit, CDN serves the content directly.
    \item If cache miss, CDN fetches content from the origin server and may cache it for future requests.
\end{itemize}

\section{Push CDN}

\subsection*{Definition}
A Push CDN is a CDN strategy where the origin system or administrator proactively uploads content to CDN edge servers before users request it.

\[
\boxed{\text{In a Push CDN, content is sent to the CDN before user demand arrives.}}
\]

\begin{itemize}
    \item The origin system decides what content should be placed on CDN servers.
    \item Useful when content is predictable and expected to be popular.
    \item Reduces first-request delay because content is already present at the edge.
    \item Requires careful management of updates, stale content, and storage usage.
\end{itemize}

\subsection*{Example}
Before a major product launch, a company may push images, videos, and static files to CDN edge servers worldwide so users can download them quickly as soon as the launch goes live.

\subsection*{When Push CDN Is Better}
\begin{itemize}
    \item Large files that many users will request.
    \item Predictable content, such as movies, software updates, or launch assets.
    \item Live event preparation where media segments or static assets need to be ready near users.
\end{itemize}

\section{Pull CDN}

\subsection*{Definition}
A Pull CDN is a CDN strategy where the CDN fetches content from the origin server only when a user requests it for the first time.

\[
\boxed{\text{In a Pull CDN, content is fetched on demand and then cached for future requests.}}
\]

\begin{itemize}
    \item The first request may be slower because the CDN must contact the origin server.
    \item Future requests are faster if the content remains cached.
    \item Useful when content demand is unpredictable.
    \item Avoids pushing unnecessary content to every edge server.
\end{itemize}

\subsection*{Example}
If an Instagram user requests an image or video that is not already cached near them, the CDN may pull it from the origin server, cache it at the edge, and serve it faster to future users in the same region.

\subsection*{When Pull CDN Is Better}
\begin{itemize}
    \item User-generated content.
    \item Long-tail content where not every item is popular.
    \item Systems where it is hard to predict which content users will request.
\end{itemize}

\subsection*{Push CDN vs Pull CDN}
\begin{itemize}
    \item Push CDN is origin-driven; content is sent before users request it.
    \item Pull CDN is request-driven; content is fetched when users request it.
    \item Push CDN is better for predictable, widely accessed content.
    \item Pull CDN is better for unpredictable or user-generated content.
\end{itemize}

\section{Time To Live (TTL)}

\subsection*{Definition}
Time To Live (TTL) is the amount of time cached content is considered valid before the CDN must recheck or refresh it from the origin server.

\[
\boxed{\text{TTL controls how long cached content can be served before checking for freshness.}}
\]

\begin{itemize}
    \item A longer TTL improves performance because the CDN contacts the origin less often.
    \item A shorter TTL improves freshness because the CDN checks for updates more frequently.
    \item TTL helps balance speed, origin load, and data correctness.
    \item If TTL expires, the CDN may revalidate or refetch the content from the origin server.
\end{itemize}

\subsection*{Example}
A website logo may have a long TTL because it rarely changes. A news homepage may have a short TTL because the content changes frequently.

\section{CDN for Live Streaming}

\subsection*{Definition}
For live streaming, a CDN distributes small video segments to edge servers in near real time instead of caching one complete static file.

\[
\boxed{\text{Live CDN delivery is fast, continuous distribution of fresh video segments to users.}}
\]

\begin{itemize}
    \item The live video feed is ingested by the streaming platform.
    \item The video is broken into small segments.
    \item CDN edge servers distribute these segments to nearby users.
    \item The goal is low buffering, high availability, and scalable delivery.
    \item Live streaming is less about long-term caching and more about fast real-time distribution.
\end{itemize}

\subsection*{Example}
For a live Netflix boxing match, the video stream is processed into small live segments and delivered through CDN edge servers near viewers. This allows millions of users to watch the event without every request hitting the main origin servers directly.

\section{CDN vs Data Center Replication}

\subsection*{Definition}
CDN caching and data center replication solve different problems. A CDN speeds up content delivery, while data center replication keeps backend data available across regions.

\[
\boxed{\text{Use CDNs for fast content delivery; use replication for backend data availability and consistency.}}
\]

\begin{itemize}
    \item CDN is mainly used for static or cacheable content such as images, videos, JavaScript, CSS, and media files.
    \item Replication is used for important backend data such as users, payments, orders, and application state.
    \item CDN content may become stale depending on TTL and cache policies.
    \item Replicated data can provide stronger consistency, but it is more expensive and complex.
    \item Real systems often use both together.
\end{itemize}

\subsection*{Example}
Netflix can replicate user account data across regions while using CDN edge servers to deliver video content. Account data needs correctness, but video delivery needs speed and scale.

\section{Load Balancer}

\subsection*{Definition}
A load balancer distributes incoming client requests across multiple backend servers to improve performance, scalability, and availability.

\[
\boxed{\text{A load balancer routes each request to an appropriate healthy server.}}
\]

\begin{itemize}
    \item Prevents a single server from receiving all traffic.
    \item Allows multiple servers to process requests in parallel.
    \item Performs health checks to avoid routing traffic to failed servers.
    \item Can operate at the transport layer or application layer.
    \item Can reroute traffic when a server or regional load balancer fails.
\end{itemize}

\subsection*{Real-World Example}
An e-commerce application has three application servers. Instead of sending every request to one server, a load balancer distributes customer requests among all three servers. If one server becomes unavailable, traffic is routed only to the remaining healthy servers.


\subsection{Layer 4 Load Balancer}

\subsection*{Definition}
A Layer 4 load balancer operates at the transport layer and routes traffic using information such as source IP address, destination IP address, and port number.

\[
\boxed{\text{L4 routing uses network and transport information without inspecting application content.}}
\]

\begin{itemize}
    \item Commonly handles TCP and UDP traffic.
    \item Does not inspect URLs, cookies, or HTTP headers.
    \item Provides fast and relatively simple traffic distribution.
    \item Can be used for HTTP, database, gaming, or other TCP/UDP connections.
\end{itemize}

\subsection*{Real-World Example}
A client creates a TCP connection to port 443. The L4 load balancer sees the destination IP and port and forwards the connection to Server A. The next TCP connection may be forwarded to Server B, regardless of the requested URL.


\subsection{Layer 7 Load Balancer}

\subsection*{Definition}
A Layer 7 load balancer operates at the application layer and routes requests by inspecting application-level information such as URLs, HTTP headers, cookies, and request methods.

\[
\boxed{\text{L7 routing uses request content to select the appropriate backend service.}}
\]

\begin{itemize}
    \item Understands application protocols such as HTTP and HTTPS.
    \item Supports path-based and header-based routing.
    \item Can route different application features to different server groups.
    \item Provides more routing flexibility but requires more processing than L4.
\end{itemize}

\subsection*{Real-World Example}
For an e-commerce application:
\begin{itemize}
    \item Requests to \texttt{/images} are routed to media servers.
    \item Requests to \texttt{/login} are routed to authentication servers.
    \item Requests to \texttt{/payments} are routed to payment servers.
\end{itemize}

\subsection*{L4 vs.\ L7 Load Balancing}

\begin{center}
\begin{tabular}{|p{3cm}|p{5cm}|p{5cm}|}
\hline
\textbf{Feature} & \textbf{Layer 4} & \textbf{Layer 7} \\ \hline
Routing information & IP address and port & URL, headers, cookies, and request content \\ \hline
Application awareness & No & Yes \\ \hline
Processing & Faster and simpler & More flexible but more processing-intensive \\ \hline
Example & Route TCP connections among servers & Route \texttt{/login} and \texttt{/images} to different services \\ \hline
\end{tabular}
\end{center}


\section{Load-Balancing Algorithms}

\subsection*{Definition}
A load-balancing algorithm determines which backend server should receive each incoming request.

\[
\boxed{\text{The algorithm defines how traffic is distributed among available servers.}}
\]

\subsection*{Round-Robin}

Round-robin sends requests to servers sequentially in a repeating order.

\begin{itemize}
    \item Request 1 goes to Server A.
    \item Request 2 goes to Server B.
    \item Request 3 goes to Server C.
    \item Request 4 returns to Server A.
    \item Works well when servers have similar capacity.
\end{itemize}

\subsection*{Least Connections}

Least connections sends a new request to the server currently handling the fewest active connections.

\begin{itemize}
    \item Useful when requests have different processing durations.
    \item Helps prevent a busy server from receiving additional traffic.
    \item If Server A has 30 active connections and Server B has 25, the next connection is sent to Server B.
\end{itemize}

\subsection*{Weighted Load Balancing}

Weighted load balancing assigns each server a weight according to its processing capacity.

\begin{itemize}
    \item More powerful servers receive a larger percentage of traffic.
    \item Useful when backend servers have different hardware or capacities.
    \item Weights can be combined with round-robin or least-connections routing.
\end{itemize}

\subsection*{Example}
Suppose Server A has weight 2 and Server B has weight 1.

\begin{itemize}
    \item Server A receives approximately two out of every three requests.
    \item Server B receives approximately one out of every three requests.
\end{itemize}

\subsection*{Geographic Routing}

Geographic routing sends users to servers located in an appropriate geographic region.

\begin{itemize}
    \item European users can be routed to European servers.
    \item Asian users can be routed to Asian servers.
    \item North American users can be routed to North American servers.
    \item Reduces long-distance network latency.
\end{itemize}


\subsection{Global and Regional Load Balancing}

\subsection*{Definition}
Global and regional load balancing uses multiple levels of load balancers to route users first to an appropriate geographic region and then to a server within that region.

\[
\boxed{\text{Global routing selects the region; regional routing selects the server.}}
\]

\begin{itemize}
    \item A global routing system directs users to a nearby or healthy region.
    \item Each region contains its own regional load balancer.
    \item The regional load balancer distributes traffic among local servers.
    \item Avoids sending every request through one distant centralized load balancer.
    \item Different regions may use different load-balancing algorithms.
\end{itemize}

\subsection*{Real-World Example}
A North American user is directed by a global routing layer to the North American region. The regional load balancer then uses least-connections routing to select one of the available application servers.

A European user is directed to the European region, where the regional load balancer may use round-robin routing because all servers have similar capacity.


\subsection{Load-Balancer Health Checks and Failover}

\subsection*{Definition}
Health checks allow a load balancer to determine whether backend servers or other load balancers are available. Failover redirects traffic when a component becomes unavailable.

\[
\boxed{\text{Health checks detect failure, and failover keeps traffic flowing through healthy components.}}
\]

\begin{itemize}
    \item The load balancer periodically checks backend server health.
    \item Unhealthy servers are temporarily removed from traffic distribution.
    \item Traffic is restored when the server becomes healthy again.
    \item A standby load balancer can replace a failed active load balancer.
    \item Regional traffic can be redirected to another region during a regional failure.
\end{itemize}

\subsection*{Real-World Example}
A regional load balancer manages three application servers. If Server B stops responding to health checks, the load balancer routes new requests only to Servers A and C.

If the entire regional load balancer fails, a secondary load balancer or another region can receive the traffic.


\subsection{Load-Balancer Replication}

\subsection*{Definition}
Load-balancer replication uses multiple load-balancer instances so that the load-balancing layer does not become a single point of failure.

\[
\boxed{\text{Multiple load balancers improve the availability of the traffic-routing layer.}}
\]

\begin{itemize}
    \item Load balancers may operate in active--active or active--passive configurations.
    \item In active--active, multiple load balancers handle traffic simultaneously.
    \item In active--passive, one load balancer handles traffic while another remains on standby.
    \item Configuration and routing information may be synchronized across instances.
\end{itemize}

\subsection*{Real-World Example}
Two load balancers serve an application in an active--active configuration. Both receive client traffic and route requests to healthy backend servers. If one load balancer fails, the other continues handling incoming requests.

\section{Monolithic Architecture}

\subsection*{Definition}
A monolithic architecture packages the application's controllers, business logic, and data-access components into one deployable application.

\begin{center}
\fbox{\parbox{0.94\linewidth}{\textbf{Quick Revision:} A monolith is developed, deployed, and scaled as one application.}}
\end{center}

\begin{itemize}
    \item All components usually run inside the same process or JVM.
    \item The application is commonly packaged as one JAR, WAR, or container.
    \item A load increase in one feature may require scaling the entire application.
    \item A failure can affect the complete application.
\end{itemize}

\subsection*{Example}
An e-commerce Spring Boot application contains user, product, order, payment, and notification service classes inside one project. If product traffic increases, another instance of the entire application must be deployed.

\section{Microservices Architecture}

\subsection*{Definition}
A microservices architecture divides a system into small, independently deployable services, where each service handles a specific business capability.

\begin{center}
\fbox{\parbox{0.94\linewidth}{\textbf{Quick Revision:} A microservice is an independently deployable and independently scalable application.}}
\end{center}

\begin{itemize}
    \item Each service runs as a separate process or container.
    \item Services communicate through HTTP, gRPC, or messaging systems.
    \item Each service can be deployed and scaled independently.
    \item Services are loosely coupled but introduce network and operational complexity.
\end{itemize}

\subsection*{Example}
An e-commerce system may contain separate Order, Payment, Inventory, and Notification microservices. During a sale, only the Order and Inventory services may require additional instances.

\subsection*{Monolith vs.\ Microservices}

\begin{center}
\begin{tabular}{|p{0.28\linewidth}|p{0.30\linewidth}|p{0.30\linewidth}|}
\hline
\textbf{Property} & \textbf{Monolith} & \textbf{Microservices} \\
\hline
Deployment & Entire application together & Each service independently \\
\hline
Scaling & Entire application & Individual service \\
\hline
Communication & Direct method calls & Network calls or messages \\
\hline
Failure isolation & Lower & Higher \\
\hline
Operational complexity & Lower & Higher \\
\hline
\end{tabular}
\end{center}

\subsection*{Spring Boot Service Classes vs.\ Microservices}

\subsection*{Definition}
A Spring Boot service class is an internal Java component containing business logic, while a microservice is a complete, separately running and deployable application.

\begin{center}
\fbox{\parbox{0.94\linewidth}{\textbf{Quick Revision:} Multiple \texttt{@Service} classes inside one Spring Boot application are not multiple microservices.}}
\end{center}

\begin{itemize}
    \item Service classes inside one application share the same JVM and deployment.
    \item They are generally called through normal Java method calls.
    \item They cannot normally be deployed or scaled independently.
    \item Spring Boot can be used to build either a monolith or individual microservices.
\end{itemize}

\subsection*{Example}
A Spring Boot project containing \texttt{UserService}, \texttt{EventService}, and \texttt{PaymentService} with one \texttt{main()} method is one monolithic application. To create independent microservices, each capability must become a separate Spring Boot application.

\subsection*{API Gateway Pattern}

\subsection*{Definition}
The API Gateway pattern provides a single entry point through which clients access multiple backend microservices.

\begin{center}
\fbox{\parbox{0.94\linewidth}{\textbf{Quick Revision:} The API Gateway receives client requests and routes them to the correct microservice.}}
\end{center}

\begin{itemize}
    \item Hides internal service locations from clients.
    \item Performs request routing, authentication, rate limiting, and logging.
    \item May combine responses from multiple services.
    \item Can become a bottleneck if not replicated and scaled properly.
\end{itemize}

\subsection*{Example}
An e-commerce client sends all requests to one gateway. The gateway routes \texttt{/users} to the User Service, \texttt{/orders} to the Order Service, and \texttt{/payments} to the Payment Service.

\subsection*{Database per Service Pattern}

\subsection*{Definition}
The Database per Service pattern gives each microservice ownership of its own data store and prevents other services from directly modifying that data.

\begin{center}
\fbox{\parbox{0.94\linewidth}{\textbf{Quick Revision:} Each microservice owns its data and exposes it through APIs or events.}}
\end{center}

\begin{itemize}
    \item Reduces database-level coupling between services.
    \item Allows each service to choose an appropriate database technology.
    \item Prevents one service from depending directly on another service's schema.
    \item Makes cross-service transactions and joins more difficult.
\end{itemize}

\subsection*{Example}
The Order Service stores order data in its database, while the Payment Service stores payment data separately. The Order Service requests payment information through the Payment Service API instead of querying its database directly.

\subsection*{Saga Pattern}

\subsection*{Definition}
The Saga pattern manages a business transaction across multiple microservices as a sequence of local transactions, with compensating transactions used to undo completed steps when a later step fails.

\begin{center}
\fbox{\parbox{0.94\linewidth}{\textbf{Quick Revision:} A Saga replaces one distributed transaction with local transactions and compensating actions.}}
\end{center}

\begin{itemize}
    \item Each microservice updates only its own database.
    \item Successful steps trigger the next step in the workflow.
    \item A failure triggers compensating operations.
    \item Compensation is a business-level undo operation, not a database rollback.
\end{itemize}

\subsection*{Example}
For an online order:
\begin{enumerate}
    \item The Order Service creates the order.
    \item The Payment Service charges the customer.
    \item The Inventory Service reserves the item.
    \item The Shipping Service creates the shipment.
\end{enumerate}

If inventory reservation fails after payment succeeds, the Payment Service refunds the customer and the Order Service cancels the order.

\subsection*{Saga Choreography vs.\ Orchestration}

\begin{center}
\begin{tabular}{|p{0.25\linewidth}|p{0.33\linewidth}|p{0.33\linewidth}|}
\hline
\textbf{Property} & \textbf{Choreography} & \textbf{Orchestration} \\
\hline
Control & Distributed among services & Central orchestrator controls workflow \\
\hline
Communication & Events & Commands and responses \\
\hline
Advantage & Loose coupling & Workflow is easier to understand \\
\hline
Limitation & Difficult to trace complex flows & Orchestrator may become complex \\
\hline
\end{tabular}
\end{center}

\subsection*{Circuit Breaker Pattern}

\subsection*{Definition}
The Circuit Breaker pattern temporarily stops requests to an unhealthy service after repeated failures, preventing resource exhaustion and cascading failures.

\begin{center}
\fbox{\parbox{0.94\linewidth}{\textbf{Quick Revision:} A circuit breaker fails fast instead of repeatedly calling an unavailable service.}}
\end{center}

\begin{itemize}
    \item \textbf{Closed:} Requests are allowed normally.
    \item \textbf{Open:} Requests fail immediately without calling the service.
    \item \textbf{Half-open:} A limited test request checks whether the service recovered.
    \item Commonly combined with timeouts, retries, and fallback responses.
\end{itemize}

\subsection*{Example}
If the Payment Service repeatedly times out, the Order Service opens its circuit breaker and temporarily rejects payment requests instead of waiting for repeated network timeouts.

\section{Service Discovery Pattern}

\subsection*{Definition}
The Service Discovery pattern allows services to locate available instances of another service without storing fixed IP addresses.

\begin{center}
\fbox{\parbox{0.94\linewidth}{\textbf{Quick Revision:} Service discovery maps a logical service name to its currently available instances.}}
\end{center}

\begin{itemize}
    \item Service instances may be created, removed, or assigned new addresses.
    \item A service registry or orchestration platform tracks healthy instances.
    \item Client-side or server-side load balancing selects an instance.
    \item Kubernetes commonly provides discovery through Services and internal DNS.
\end{itemize}

\subsection*{Example}
The Order Service calls \texttt{payment-service} rather than a fixed server IP. Service discovery resolves the name to one of the available Payment Service instances.

\section{Event-Driven Communication Pattern}

\subsection*{Definition}
Event-driven communication allows a service to publish an event that can be consumed asynchronously by one or more interested services.

\begin{center}
\fbox{\parbox{0.94\linewidth}{\textbf{Quick Revision:} Publishers announce what happened without directly calling every consumer.}}
\end{center}

\begin{itemize}
    \item Producers and consumers are loosely coupled.
    \item Multiple services can react to the same event.
    \item Message brokers such as Kafka can store and distribute events.
    \item Consumers must handle duplicate, delayed, or out-of-order events.
\end{itemize}

\subsection*{Example}
After completing an order, the Order Service publishes an \texttt{OrderCompleted} event. Notification, Analytics, and Loyalty services consume the event independently.

\section{Timeout and Retry Patterns}

\subsection*{Definition}
A timeout limits how long a service waits for a response, while a retry repeats a failed operation when the failure may be temporary.

\begin{center}
\fbox{\parbox{0.94\linewidth}{\textbf{Quick Revision:} Timeouts prevent indefinite waiting, while limited retries handle temporary failures.}}
\end{center}

\begin{itemize}
    \item Every network request should have a reasonable timeout.
    \item Retries should be limited and use increasing delays.
    \item Excessive retries can increase load during an outage.
    \item Retried operations should be idempotent whenever possible.
\end{itemize}

\subsection*{Example}
The Order Service waits two seconds for the Payment Service. If the request fails temporarily, it retries twice with increasing delays before returning an error.

\subsection*{Bulkhead Pattern}

\subsection*{Definition}
The Bulkhead pattern isolates resources used by different operations so that failure or overload in one area does not exhaust resources required by the entire application.

\begin{center}
\fbox{\parbox{0.94\linewidth}{\textbf{Quick Revision:} Bulkheads limit the damage caused by one slow or failing dependency.}}
\end{center}

\begin{itemize}
    \item Separates thread pools, connection pools, or concurrency limits.
    \item Prevents one dependency from consuming all application resources.
    \item Improves fault isolation and availability.
    \item Resource limits must be configured according to expected traffic.
\end{itemize}

\subsection*{Example}
An application reserves separate thread pools for Payment and Product requests. If Payment requests become slow, Product requests can still be processed.

\section{CQRS Pattern}

\subsection*{Definition}
Command Query Responsibility Segregation separates operations that modify data from operations that read data.

\begin{center}
\fbox{\parbox{0.94\linewidth}{\textbf{Quick Revision:} CQRS uses separate models or paths for writes and reads.}}
\end{center}

\begin{itemize}
    \item Commands create, update, or delete data.
    \item Queries retrieve data without modifying it.
    \item Read and write workloads can be optimized and scaled separately.
    \item Separate read models may temporarily lag behind the write model.
\end{itemize}

\subsection*{Example}
An order system sends order creation and cancellation commands to a write database, while customer order-history requests are served from a separate read-optimized database.

\section{Strangler Fig Pattern}

\subsection*{Definition}
The Strangler Fig pattern gradually replaces parts of a monolithic application with new services instead of rewriting the entire system at once.

\begin{center}
\fbox{\parbox{0.94\linewidth}{\textbf{Quick Revision:} Extract monolithic functionality incrementally until the old implementation can be removed.}}
\end{center}

\begin{itemize}
    \item New requests are gradually routed to extracted services.
    \item Unchanged functionality continues running in the monolith.
    \item Reduces the risk of a complete system rewrite.
    \item Requires temporary integration between the monolith and new services.
\end{itemize}

\subsection*{Example}
A company first extracts payment processing from its e-commerce monolith into a Payment microservice. It later extracts Inventory and Order services while the remaining features continue running in the monolith.

\section{Relational and NoSQL Databases}

\subsection*{Definition}
Relational databases store structured data in tables with predefined schemas, while NoSQL databases use flexible models such as key-value, document, wide-column, and graph storage.

\begin{center}
\fbox{\parbox{0.94\linewidth}{\textbf{Revision:} SQL databases emphasize relationships, constraints, and transactions; NoSQL databases emphasize flexible schemas and distributed scalability.}}
\end{center}

\begin{itemize}
\item SQL databases support tables, joins, foreign keys, constraints, and ACID transactions.
\item Examples include PostgreSQL, MySQL, Oracle, and SQL Server.
\item NoSQL databases are commonly selected for flexible schemas, high-volume workloads, or horizontal scaling.
\item NoSQL does not mean that SQL-like querying is always unavailable.
\item Large systems commonly use multiple database types for different workloads.
\end{itemize}

\subsection*{Example}
An e-commerce platform may use PostgreSQL for orders and payments, MongoDB for the product catalog, Redis for shopping carts, Cassandra for activity events, and a graph database for recommendations.

\section{ACID Properties}

\subsection*{Definition}
ACID properties ensure that database transactions execute reliably while preserving application-defined constraints during failures and concurrent operations.

\begin{center}
\fbox{\parbox{0.94\linewidth}{\textbf{Revision:} ACID stands for Atomicity, Consistency, Isolation, and Durability.}}
\end{center}

\begin{itemize}
\item \textbf{Atomicity:} Either every operation in a transaction succeeds or the entire transaction is rolled back.
\item \textbf{Consistency:} A transaction preserves database constraints and invariants.
\item \textbf{Isolation:} Concurrent transactions behave according to a defined isolation level and do not improperly interfere.
\item \textbf{Durability:} Once committed, data remains stored even after a crash.
\end{itemize}

\subsection*{Example}
During a bank transfer, deducting money from one account and crediting another must occur as one transaction. A partial transfer must never be committed.

\subsection*{Important Distinction}
ACID consistency refers to preserving database rules. CAP consistency refers to distributed replicas presenting a single up-to-date view of data.

\section{Database Normalization}

\subsection*{Definition}
Normalization organizes relational data into separate tables to reduce unnecessary duplication and prevent update, insertion, and deletion anomalies.

\begin{center}
\fbox{\parbox{0.94\linewidth}{\textbf{Revision:} Normalization reduces redundancy by storing each fact in an appropriate table.}}
\end{center}

\begin{itemize}
\item \textbf{First Normal Form (1NF):} Each column contains an atomic value, with no repeating groups.
\item \textbf{Second Normal Form (2NF):} The table is in 1NF and every non-key attribute depends on the complete composite key.
\item \textbf{Third Normal Form (3NF):} The table is in 2NF and contains no transitive dependency between non-key attributes.
\item BCNF, 4NF, and 5NF address more advanced dependency cases.
\end{itemize}

\subsection*{Example}
Instead of repeatedly storing a department phone number in every student row, the database stores departments separately and references them through a department identifier.

\begin{verbatim}
Student(student_id, name, department_id)
Department(department_id, department_name, phone)
\end{verbatim}

\section{Key-Value Databases}

\subsection*{Definition}
A key-value database stores each value under a unique key and retrieves it primarily through direct key lookup.

\begin{center}
\fbox{\parbox{0.94\linewidth}{\textbf{Revision:} A key-value store behaves like a distributed dictionary optimized for fast lookups.}}
\end{center}

\begin{itemize}
\item The key uniquely identifies the stored value.
\item Values may contain strings, numbers, lists, sets, or serialized objects.
\item Queries are generally based on keys rather than arbitrary internal fields.
\item Common examples include Redis and Amazon DynamoDB.
\end{itemize}

\subsection*{Example}
An e-commerce application can store a shopping cart using the user identifier as the key.

\begin{verbatim}
Key: cart:user:101
Value: {items: [901, 902], total_items: 2}
\end{verbatim}

\section{Document Databases}

\subsection*{Definition}
A document database stores records as self-contained JSON-like documents whose fields may be nested and may vary between documents.

\begin{center}
\fbox{\parbox{0.94\linewidth}{\textbf{Revision:} Document databases store flexible objects and allow queries on fields inside those objects.}}
\end{center}

\begin{itemize}
\item Documents are grouped into collections.
\item Documents in the same collection do not require identical fields.
\item Nested objects and arrays can be stored directly.
\item Fields can be queried, projected, sorted, aggregated, and indexed.
\item MongoDB is a document database, not merely a key-value database.
\end{itemize}

\subsection*{Example}
A MongoDB product collection can store different attributes for different product categories.

\begin{verbatim}
{
"_id": 901,
"name": "Laptop",
"specifications": {
"ram": "16 GB",
"storage": "1 TB"
},
"tags": ["electronics", "computer"]
}
\end{verbatim}

\subsection*{MongoDB Query Example}

\begin{verbatim}
db.products.find(
{"specifications.ram": "16 GB"},
{"name": 1, "specifications": 1, "_id": 0}
)
\end{verbatim}

\section{Wide-Column Databases}

\subsection*{Definition}
A wide-column database stores data using partition keys, rows, and flexible groups of columns optimized for distributed storage and high write throughput.

\begin{center}
\fbox{\parbox{0.94\linewidth}{\textbf{Revision:} Wide-column databases organize data around partition keys and query-specific access patterns.}}
\end{center}

\begin{itemize}
\item Related rows are grouped using a partition key.
\item Different rows may contain different columns.
\item Data models are designed around known queries rather than joins.
\item Common examples include Cassandra and HBase.
\item They are suitable for large-scale event, logging, and time-series workloads.
\end{itemize}

\subsection*{Example}
A Cassandra table may store sensor readings by device and timestamp.

\begin{verbatim}
device_id | timestamp | temperature | battery
---------------------------------------------

A101      | 10:00     | 24.5        | 91
A101      | 10:05     | 25.0        | 90
\end{verbatim}

All readings for the same device can be stored together and retrieved efficiently.

\section{Graph Databases}

\subsection*{Definition}
A graph database represents entities as nodes and relationships between entities as edges.

\begin{center}
\fbox{\parbox{0.94\linewidth}{\textbf{Revision:} Graph databases are optimized for traversing relationships rather than repeatedly joining tables.}}
\end{center}

\begin{itemize}
\item Nodes represent entities such as users, products, or companies.
\item Edges represent relationships such as \texttt{FRIEND_OF}, \texttt{BOUGHT}, or \texttt{WORKS_AT}.
\item Nodes and edges may contain properties.
\item Common use cases include fraud detection, social networks, knowledge graphs, and recommendations.
\item Neo4j is a common graph database.
\end{itemize}

\subsection*{Example}

\begin{verbatim}
(Alice)-[:BOUGHT]->(Laptop)
(Bob)-[:BOUGHT]->(Laptop)
(Bob)-[:BOUGHT]->(Mouse)
\end{verbatim}

The system can recommend the mouse to Alice by traversing purchases made by users with similar behavior.

\section{Database Replication}

\subsection*{Definition}
Replication maintains copies of the same database records on multiple nodes to improve availability, fault tolerance, geographic access, and read scalability.

\begin{center}
\fbox{\parbox{0.94\linewidth}{\textbf{Revision:} Replication copies the same data across nodes; it does not divide the dataset.}}
\end{center}

\begin{itemize}
\item In primary-replica replication, clients send writes to the primary node.
\item The primary propagates committed changes to one or more replicas.
\item Replicas may serve read requests, although asynchronous replicas can return stale data.
\item If the primary fails, an eligible replica can be promoted.
\item In multi-primary replication, multiple nodes accept writes and conflicts must be resolved.
\item Replication is not a replacement for backups because accidental deletions can also be replicated.
\end{itemize}

\subsection*{Synchronous and Asynchronous Replication}

\begin{itemize}
\item \textbf{Synchronous replication:} The primary waits for replica acknowledgement before confirming the write.
\item \textbf{Asynchronous replication:} The primary confirms the write before every replica receives it.
\item Synchronous replication reduces data-loss risk but increases latency.
\item Asynchronous replication improves performance but introduces replication lag.
\end{itemize}

\subsection*{Example}
A PostgreSQL order database may have one primary and two read replicas. Writes go to the primary, while order-history queries can be distributed across the replicas.

\section{Database Sharding}

\subsection*{Definition}
Sharding horizontally partitions a large dataset and stores different subsets of the data on different database servers.

\begin{center}
\fbox{\parbox{0.94\linewidth}{\textbf{Revision:} Sharding distributes different data across nodes to scale storage and traffic horizontally.}}
\end{center}

\begin{itemize}
\item A shard key determines where each record is stored.
\item Queries containing the shard key can be routed directly to the correct shard.
\item Queries without the shard key may contact every shard.
\item Cross-shard joins and transactions are more complex.
\item Poor shard keys can create uneven load or hot shards.
\item Each shard can also be replicated for availability.
\end{itemize}

\subsection*{Example}
An order system can shard records using a hash of \texttt{customer_id}.

\begin{verbatim}
shard_number = hash(customer_id) mod number_of_shards
\end{verbatim}

All orders for customer 101 can be routed to the shard calculated from that customer's identifier.

\subsection*{Replication versus Sharding}

\begin{itemize}
\item Replication stores copies of the same data.
\item Sharding stores different portions of the data.
\item Replication mainly improves availability and read capacity.
\item Sharding mainly improves storage capacity and distributed throughput.
\end{itemize}

\section{Database Federation}

\subsection*{Definition}
Database federation, also called functional partitioning, separates data into independent databases according to business domains.

\begin{center}
\fbox{\parbox{0.94\linewidth}{\textbf{Revision:} Federation separates databases by business responsibility, while sharding divides records within one large dataset.}}
\end{center}

\begin{itemize}
\item Related tables are grouped by domain rather than creating one database for every table.
\item Each domain can scale and deploy independently.
\item Failures can be isolated to one business function.
\item Cross-database joins and transactions become more difficult.
\item Different domains may use different database technologies.
\end{itemize}

\subsection*{Example}
An e-commerce platform may use separate databases for customers, products, orders, and payments.

\begin{verbatim}
Customer Database: customers, addresses, preferences
Order Database: orders, order_items, shipments
Payment Database: payments, refunds, payment_attempts
\end{verbatim}

\section{Database Denormalization}

\subsection*{Definition}
Denormalization intentionally duplicates or embeds selected data to reduce joins and improve read performance.

\begin{center}
\fbox{\parbox{0.94\linewidth}{\textbf{Revision:} Denormalization exchanges additional storage and update complexity for faster reads.}}
\end{center}

\begin{itemize}
\item Frequently accessed data can be copied into the same row or document.
\item Reads become simpler because fewer joins or network calls are required.
\item Updates become harder because duplicated values may need synchronization.
\item Denormalization should follow actual query patterns.
\item Historical snapshots are a useful form of intentional duplication.
\end{itemize}

\subsection*{SQL Example}
A normalized order query may join \texttt{orders}, \texttt{customers}, \texttt{order_items}, and \texttt{products}. A denormalized order can store the customer name and product name at purchase time.

\begin{verbatim}
orders(
order_id,
customer_id,
customer_name_at_purchase,
total
)

order_items(
order_id,
product_id,
product_name_at_purchase,
unit_price_at_purchase
)
\end{verbatim}

\subsection*{MongoDB Example}
MongoDB can embed order items and customer snapshot data inside one order document.

\begin{verbatim}
{
"_id": 5001,
"customerSnapshot": {
"customerId": 101,
"name": "Alice"
},
"items": [
{
"productId": 901,
"productName": "Laptop",
"unitPrice": 1000
}
]
}
\end{verbatim}

\section{Database Indexing}

\subsection*{Definition}
An index is an additional data structure that allows a database to locate matching rows or documents without scanning the complete dataset.

\begin{center}
\fbox{\parbox{0.94\linewidth}{\textbf{Revision:} Indexes speed up reads but consume storage and add overhead to writes.}}
\end{center}

\begin{itemize}
\item Indexes are useful for frequently filtered, joined, sorted, or uniquely constrained fields.
\item A full table or collection scan examines every record.
\item A composite index contains multiple fields, and field order matters.
\item Excessive indexes slow inserts, updates, and deletes.
\item Indexes should be created according to real query patterns.
\end{itemize}

\subsection*{SQL Example}

\begin{verbatim}
CREATE UNIQUE INDEX idx_customers_email
ON customers(email);

CREATE INDEX idx_orders_customer_created
ON orders(customer_id, created_at DESC);
\end{verbatim}

The second index supports retrieving a customer's newest orders efficiently.

\subsection*{MongoDB Example}

\begin{verbatim}
db.customers.createIndex(
{email: 1},
{unique: true}
)

db.orders.createIndex({
customerId: 1,
createdAt: -1
})
\end{verbatim}

\subsection*{Important Distinction}
A primary key identifies a record and normally creates an index automatically. An ordinary index improves access speed but does not necessarily enforce uniqueness.

\section{SQL Query Tuning}

\subsection*{Definition}
SQL query tuning improves database performance by examining execution plans and optimizing queries, indexes, joins, schema design, and data access patterns.

\begin{center}
\fbox{\parbox{0.94\linewidth}{\textbf{Revision:} SQL tuning begins with measuring the execution plan rather than guessing which query is slow.}}
\end{center}

\begin{itemize}
\item Use \texttt{EXPLAIN} or \texttt{EXPLAIN ANALYZE} to inspect the execution plan.
\item Add indexes for important filters, joins, and sorting operations.
\item Retrieve only required columns instead of using \texttt{SELECT *}.
\item Filter rows as early as possible.
\item Avoid unnecessary joins and repeated calculations.
\item Use suitable pagination, caching, summary tables, or materialized views where appropriate.
\item Keep database statistics updated so the optimizer can estimate row counts accurately.
\end{itemize}

\subsection*{Example}

\begin{verbatim}
EXPLAIN ANALYZE
SELECT
o.order_id,
o.status,
o.total,
o.created_at
FROM customers AS c
JOIN orders AS o
ON o.customer_id = c.customer_id
WHERE c.email = '[alice@example.com](mailto:alice@example.com)'
ORDER BY o.created_at DESC
LIMIT 20;
\end{verbatim}

Useful supporting indexes are:

\begin{verbatim}
CREATE UNIQUE INDEX idx_customers_email
ON customers(email);

CREATE INDEX idx_orders_customer_created
ON orders(customer_id, created_at DESC);
\end{verbatim}

\section{Join Order and Early Filtering}

\subsection*{Definition}
Join optimization reduces the size of intermediate results by filtering selective data before joining it with much larger tables.

\begin{center}
\fbox{\parbox{0.94\linewidth}{\textbf{Revision:} Reduce the dataset before performing an expensive join whenever the execution plan permits it.}}
\end{center}

\begin{itemize}
\item Joining large tables before applying selective filters can produce huge intermediate results.
\item Joining a small filtered result with larger tables usually requires less work.
\item Modern SQL optimizers can reorder inner joins automatically.
\item Changing the written join order alone may not change the execution plan.
\item The actual plan must be verified using \texttt{EXPLAIN ANALYZE}.
\end{itemize}

\subsection*{Expensive Query}

\begin{verbatim}
SELECT
o.order_id,
p.product_name,
oi.quantity
FROM order_items AS oi
JOIN products AS p
ON p.product_id = oi.product_id
JOIN orders AS o
ON o.order_id = oi.order_id
WHERE o.customer_id = 101;
\end{verbatim}

This plan can be expensive if a large \texttt{order_items}--\texttt{products} result is created before filtering the customer's orders.

\subsection*{Improved Query}

\begin{verbatim}
SELECT
filtered_orders.order_id,
p.product_name,
oi.quantity
FROM (
SELECT order_id
FROM orders
WHERE customer_id = 101
) AS filtered_orders
JOIN order_items AS oi
ON oi.order_id = filtered_orders.order_id
JOIN products AS p
ON p.product_id = oi.product_id;
\end{verbatim}

The query first identifies a small number of orders and then joins only their order items and products.

\subsection*{Supporting Indexes}

\begin{verbatim}
CREATE INDEX idx_orders_customer
ON orders(customer_id, order_id);

CREATE INDEX idx_order_items_order
ON order_items(order_id);
\end{verbatim}

\section{Cache}

\subsection*{Definition}
A cache is a temporary, high-speed storage layer that keeps copies of frequently or recently accessed data closer to the requester. It reduces retrieval latency and decreases repeated load on the primary data source.

\begin{center}
\fbox{\parbox{0.92\linewidth}{\textbf{Revision:} A cache improves performance by serving reusable data from faster storage instead of repeatedly accessing the source of truth.}}
\end{center}

\begin{itemize}
\item Cached data is usually a copy of data stored in a database, API, file system, or origin server.
\item The primary database or origin server normally remains the source of truth.
\item Caching mainly improves latency and reduces backend load.
\item It may improve availability when cached data can be served during a temporary backend failure.
\item Cached data may become stale if the source changes before the cache is updated.
\end{itemize}

\subsection*{Real-World Example}
An e-commerce application stores frequently requested product details in Redis. When another user opens the same product page, the application retrieves the details from Redis instead of querying the database again.

\section{Time-to-Live (TTL)}

\subsection*{Definition}
Time-to-Live is the duration for which a cached entry remains valid. After its TTL expires, the entry is considered invalid and must normally be removed or refreshed before being served again.

\begin{center}
\fbox{\parbox{0.92\linewidth}{\textbf{Revision:} TTL controls how long cached data may be reused before it must be refreshed from the source of truth.}}
\end{center}

\begin{itemize}
\item A short TTL provides fresher data but causes more database or origin requests.
\item A long TTL reduces backend load but increases the possibility of stale data.
\item An expired entry may be deleted immediately, during the next access, or through background cleanup.
\item The first request after expiration may cause a cache miss.
\item Many popular entries expiring together can create a sudden increase in backend traffic.
\end{itemize}

\subsection*{Real-World Example}
A weather API response is cached for five minutes. Requests arriving during those five minutes use the cached response, while the first request after expiration fetches updated weather data from the API.

\section{Refresh-Ahead Caching}

\subsection*{Definition}
Refresh-ahead caching proactively reloads frequently accessed data from the source of truth shortly before its TTL expires.

\begin{center}
\fbox{\parbox{0.92\linewidth}{\textbf{Revision:} Refresh-ahead updates popular cache entries before expiration to reduce cache misses and backend load.}}
\end{center}

\begin{itemize}
\item The system identifies frequently accessed or important entries.
\item A background process refreshes an entry before its expiration time.
\item Users can continue receiving data from the cache without waiting for a database query.
\item It reduces latency spikes caused by popular entries expiring.
\item It may waste resources if refreshed data is not requested again.
\item The cache can still temporarily contain stale data if the source changes between refreshes.
\end{itemize}

\subsection*{Real-World Example}
A news application refreshes its homepage headlines in the cache every few minutes before their TTL expires. Users continue receiving the cached homepage while the refresh occurs in the background.

\section{Write-Behind Caching}

\subsection*{Definition}
Write-behind caching writes data to the cache first and updates the primary database asynchronously through a background process, queue, or worker.

\begin{center}
\fbox{\parbox{0.92\linewidth}{\textbf{Revision:} Write-behind provides fast writes by delaying database persistence, but introduces consistency and data-loss risks.}}
\end{center}

\begin{itemize}
\item The application writes the new value to the cache.
\item Success may be returned before the database is updated.
\item A background worker later persists the change to the database.
\item Subsequent reads can immediately receive the latest cached value.
\item Write latency is lower because the application does not wait for the database.
\item The cache and database remain temporarily inconsistent.
\item Retries, ordering, duplicate writes, and failure recovery must be handled carefully.
\end{itemize}

\subsection*{Data-Loss Risk}
If the cache or background worker fails before the asynchronous database write completes, an acknowledged update may be lost. Durable queues, replication, and persistent logs can reduce this risk.

\subsection*{Real-World Example}
An analytics system records view counts in Redis and periodically writes aggregated counts to the database. This supports a high write rate, but recent counts could be lost if Redis fails before they are persisted.

\section{Write-Through Caching}

\subsection*{Definition}
Write-through caching synchronously updates the primary database when data is written through the cache. The write is acknowledged only after the database update succeeds.

\begin{center}
\fbox{\parbox{0.92\linewidth}{\textbf{Revision:} Write-through keeps the cache and database synchronized by persisting every write before returning success.}}
\end{center}

\begin{itemize}
\item The application sends the write through the cache layer.
\item The cache synchronously writes the value to the database.
\item The cache stores the updated value after a successful database write.
\item Future reads can retrieve the recently written value from the cache.
\item Write latency is higher because the client waits for database persistence.
\item The risk of losing an acknowledged write is lower than with write-behind.
\end{itemize}

\subsection*{Important Distinction}
Write-through describes the write path. Read-through caching describes a read path in which the cache itself fetches missing data from the database.

\subsection*{Real-World Example}
A banking application updates a user's account preferences through a write-through cache. The cache does not confirm the update until the preferences are successfully stored in the database.

\section{Cache-Aside Caching}

\subsection*{Definition}
Cache-aside is a caching strategy in which the application explicitly manages cache reads, cache misses, and cache invalidation.

\begin{center}
\fbox{\parbox{0.92\linewidth}{\textbf{Revision:} In cache-aside, the application checks the cache first and loads data from the database only when necessary.}}
\end{center}

\subsection*{Read Flow}
\begin{itemize}
\item The application first checks the cache.
\item On a cache hit, it returns the cached value.
\item On a cache miss, it queries the database.
\item The application stores the database result in the cache.
\item The result is returned to the user.
\end{itemize}

\subsection*{Write Flow}
\begin{itemize}
\item The application updates the database.
\item It deletes or invalidates the corresponding cache entry.
\item The next read reloads the updated value into the cache.
\end{itemize}

\subsection*{Limitations}
\begin{itemize}
\item The first request after expiration has higher latency.
\item Stale data may be served if invalidation is delayed or fails.
\item Many simultaneous cache misses can create a cache stampede.
\item The application must contain additional cache-management logic.
\end{itemize}

\subsection*{Real-World Example}
A product service checks Redis for product information. If the product is missing, the service reads it from PostgreSQL, stores it in Redis, and returns it to the user. When the product changes, the service updates PostgreSQL and removes the old Redis entry.

\section{Caching Locations and Use Cases}

\subsection*{Definition}
Caching can be implemented at different layers of a system depending on where data is requested, processed, and reused.

\begin{center}
\fbox{\parbox{0.92\linewidth}{\textbf{Revision:} Caches may exist on the client, CDN, web server, application layer, or inside the database.}}
\end{center}

\subsection*{Client-Side Cache}
\begin{itemize}
\item Stored on the user's browser, mobile device, or local application memory.
\item Commonly stores images, JavaScript, CSS, API responses, and application data.
\item Reduces network requests and improves user-perceived latency.
\end{itemize}

\subsection*{CDN Cache}
\begin{itemize}
\item Stores content on geographically distributed edge servers.
\item Serves users from a nearby location instead of repeatedly contacting the origin server.
\item Commonly caches images, videos, scripts, stylesheets, and downloadable files.
\end{itemize}

\subsection*{Web-Server Cache}
\begin{itemize}
\item Stores complete pages, rendered templates, static files, or API responses on the server side.
\item May be implemented using tools such as Nginx or Varnish.
\item A single cached response can be shared across many users.
\end{itemize}

\subsection*{Application Cache}
\begin{itemize}
\item Exists inside or near the backend application.
\item May use local memory, Redis, or Memcached.
\item Commonly stores sessions, product details, permissions, configuration, and expensive computation results.
\end{itemize}

\subsection*{Database Cache}
\begin{itemize}
\item A database may internally cache frequently accessed data pages, index pages, and query plans in memory.
\item This reduces repeated disk access.
\item An external Redis cache placed before a database is normally considered an application or distributed cache.
\end{itemize}

\subsection*{Real-World Example}
When a user opens an e-commerce product page, the browser may reuse cached images, the CDN may serve static files, the web server may return cached HTML, the application may retrieve product data from Redis, and the database may use its internal buffer cache for any remaining query.


\end{document}
